% Generated draft: requires mathematical and source review.
\documentclass[11pt]{article}
\usepackage[margin=25mm]{geometry}
\usepackage{fontspec}
\setmainfont{Latin Modern Roman}
\usepackage{amsmath,amssymb,mathtools,unicode-math}
\setmathfont{Latin Modern Math}
\usepackage{enumitem,needspace,xurl,hyperref}
\setlength{\parindent}{0pt}
\setlength{\parskip}{5pt}
\setlength{\emergencystretch}{3em}
\newcommand{\N}{\mathbb{N}}
\newcommand{\Z}{\mathbb{Z}}
\newcommand{\Q}{\mathbb{Q}}
\newcommand{\R}{\mathbb{R}}
\newcommand{\C}{\mathbb{C}}
\newcommand{\F}{\mathbb{F}}
\newcommand{\A}{\mathbb{A}}
\newcommand{\PP}{\mathbb{P}}
\newcommand{\E}{\mathbb{E}}
\newcommand{\eps}{\varepsilon}
\newcommand{\dd}{\,\mathrm{d}}
\begin{document}
\section*{UnsolvedMath ID: 701410}
\textbf{Source problem code:} catalogue entry
\subsection{Definitions and mathematical statement}
Let \(r\geq 1\), let \(F=F(x_1,\ldots,x_r)\) be the free group on the variables of a word \(w\), and assume \(w\neq 1\) in \(F\). For a group \(G\), define
\[
V_w(G)=\{w(g_1,\ldots,g_r):g_1,\ldots,g_r\in G\},
\qquad w(G)=\langle V_w(G)\rangle.
\]
The width of \(w\) in \(G\) is at most \(n\) if every element of \(w(G)\) is a product of at most \(n\) elements of \(V_w(G)\cup V_w(G)^{-1}\). Since \(w(1,\ldots,1)=1\), this also permits a product of exactly \(n\) such factors by padding with identity values.

A word \(w\) is \emph{uniformly elliptic} in a class \(\mathcal C\) if there exists a function \(f:\N_{\geq 1}\to\N_{\geq 1}\) such that, for every \(d\geq 1\) and every group \(G\in\mathcal C\) generated by at most \(d\) elements, the width of \(w\) in \(G\) is at most \(f(d)\). The function may depend on \(w\) and \(\mathcal C\), but not on \(G\).

Write \(F'=[F,F]\), \(F''=[F',F']\), and, for each prime \(p\),
\[
(F')^p=\langle u^p:u\in F'\rangle.
\]
Thus \(F''(F')^p\) denotes the product of these normal subgroups of \(F\).

Let \(\mathcal C_{\mathrm{sol}}\) be the class of all finite soluble groups and let \(\mathcal C_{\mathrm{fin}}\) be the class of all finite groups. Here a group is soluble if its derived series reaches the trivial subgroup after finitely many steps.

\textbf{Problem (D. Segal).} Determine whether each of the following two assertions holds: for every nonidentity word \(w\in F\),
\[
\begin{aligned}
\text{(a)}\quad &w\text{ is uniformly elliptic in }\mathcal C_{\mathrm{sol}}
\quad\Longleftrightarrow\quad
(\forall p\text{ prime})\;w\notin F''(F')^p;\\
\text{(b)}\quad &w\text{ is uniformly elliptic in }\mathcal C_{\mathrm{fin}}
\quad\Longleftrightarrow\quad
(\forall p\text{ prime})\;w\notin F''(F')^p.
\end{aligned}
\]

For background, the supplied record attributes to A. Jaikin-Zapirain the result that, for each fixed prime \(p\), a nonidentity word \(w\) is uniformly elliptic in the class of finite \(p\)-groups if and only if \(w\notin F''(F')^p\).

The record also states that, for a class \(\mathcal C\) of finite groups, uniform ellipticity is equivalent to the algebraic verbal subgroup \(w(G)\) being closed in every finitely generated pro-\(\mathcal C\) group \(G\). A pro-\(\mathcal C\) group means an inverse limit of finite groups in \(\mathcal C\), equipped with its inverse-limit topology; finite generation here means topological finite generation. The record supplies no closure hypotheses on a general class \(\mathcal C\), so its unrestricted formulation is retained as a source qualification requiring review. The two classes in parts (a) and (b) are explicitly specified above.
\subsection{Short English statement}
For each nonidentity group word, ask whether its width admits a bound depending only on the number of generators throughout all finite soluble groups, and separately throughout all finite groups, exactly when the word lies outside \(F''(F')^p\) for every prime \(p\). The supplied record attributes the corresponding fixed-prime result for finite \(p\)-groups to A. Jaikin-Zapirain.
\subsection{Sources}
\begin{itemize}
\item Local UnsolvedMath registry, permanent ID \(701410\), exact source problem code \texttt{catalogue entry}, title \emph{Kourovka Notebook Problem 17.109}. The supplied record attributes the problem to D. Segal, with proposed year 2010, and labels it \emph{open}. This label is supplied registry data, not an independently verified status assessment. The record was created and updated on \texttt{2026-10-16T14:00:00Z}.
\item \emph{The Kourovka Notebook (Unsolved Problems in Group Theory)}, Problem 17.109, 17th issue (2010); the supplied background also identifies the 21st issue edition (2026), \texttt{arXiv:1401.0300v48}. The registry describes the Issues 1--21 collection as edited by E. I. Khukhro and V. D. Mazurov.
\item A. Jaikin-Zapirain, \emph{Revista Mat. Iberoamericana} \textbf{24} (2008), 617--630. These are the bibliographic details supplied for the finite \(p\)-group characterization; no article title is supplied.
\item D. Segal, \emph{Words: notes on verbal width in groups}, London Mathematical Society Lecture Note Series \textbf{361}, Cambridge University Press, 2009, Chapter 4, cited as further background in the supplied statement.
\end{itemize}
\end{document}
